\documentclass[12pt, a4paper]{article}

\usepackage[T1]{fontenc}
\usepackage[utf8]{inputenc}
\usepackage[british]{babel}
\usepackage{mathtools}
\usepackage{amsthm}
\usepackage{libertine}
\usepackage[libertine]{newtxmath}
\usepackage[mathscr]{euscript}
\usepackage{listings}
\usepackage{enumitem}
\usepackage{tikz-cd}
\usetikzlibrary{decorations.markings}
\usepackage{float}
\usepackage{geometry}
 \geometry{
 a4paper,
 total={170mm,257mm},
 left=20mm,
 top=30mm,
 }
\usepackage{fancyhdr}
\usepackage{hyperref}
\urlstyle{same}
\usepackage[noabbrev]{cleveref}
\usepackage{xcolor}
\usepackage{tcolorbox}
\tcbuselibrary{listings}

% Colors to match VS Code Lean 4 Light Theme
\definecolor{keywordcolor}{rgb}{0.0, 0.0, 1.0}   % blue
\definecolor{tacticcolor}{rgb}{0.0, 0.0, 1.0}    % blue
\definecolor{commentcolor}{rgb}{0.0, 0.5, 0.0}   % green
\definecolor{symbolcolor}{rgb}{0.0, 0.0, 0.0}    % black
\definecolor{sortcolor}{rgb}{0.0, 0.0, 1.0}      % blue
\definecolor{attributecolor}{rgb}{0.0, 0.0, 0.0} % black
% Light gray background color
\definecolor{codebg}{rgb}{0.97, 0.97, 0.97}

\def\lstlanguagefiles{lstlean.tex}

\newtcblisting{leancode}{
  listing only,
  listing options={
    language=lean,
    emph={sorry},
    emphstyle=\color{red},
    basicstyle=\ttfamily\footnotesize
  },
  colback=codebg,
  colframe=codebg,
  boxrule=0pt,
  arc=3mm,
  auto outer arc,
  top=10pt,
  bottom=10pt,
  left=5pt,
  right=5pt,
  width=\textwidth,
}

\theoremstyle{definition}
\newtheorem{innercustomthm}{Exercise}
\newenvironment{xca}[1] {\renewcommand\theinnercustomthm{#1}\innercustomthm}
  {\endinnercustomthm}

\pagestyle{fancy}
\fancyhf{}
\rhead{Lean Algebra Exercises{\renewcommand{\thefootnote}{\fnsymbol{footnote}}\footnotemark[2]}}
\chead{Equivalence Relation}
\lhead{Pedro N\'{u}\~{n}ez}
% \rfoot{Page \thepage}

\title{Lean Algebra Exercise 3.2}
\author{Pedro N\'{u}\~{n}ez}

\setcounter{tocdepth}{1}
\sloppy
\makeatletter
\hypersetup{
  pdfauthor={\@author},
  pdftitle={\@title},
  colorlinks,
  linkcolor=[rgb]{0.2,0.2,0.6},
  citecolor=[rgb]{0.2,0.2,0.6},
  urlcolor=[rgb]{0.2,0.2,0.6}}
\makeatother

\begin{document}

\begin{xca}{3.2}
  Let $G$ be a group and let $H \subset G$ be a subgroup.
  \begin{enumerate}[label=(\alph*)]
    \item Show that $g_{1} \sim g_{2} : \Leftrightarrow g_{1}^{-1}g_{2} \in H$ defines an equivalence relation on $G$.
    \item Use part (a) of this exercise to show that any two cosets in $G$ are either equal or disjoint.
  \end{enumerate}

  You can try to \href{https://live.lean-lang.org/#codez=JYWwDg9gTgLgBAWQIYwBYBtgCMB0BBdAcwFMsokcBxKCAVzBwGVatCb6cAhJAZ2AGMAUIIBuSKMCRZ0xOAG9KcAFxwAKgE8wxAL5wA2tTpg4lALpwAFAAllcZq3bHKASmEATYgDM4AJWLpLGxV7NiMTZ0tCOEIAclsXWwAFGmMlAF5owG8CAE64ACpouMAIIjgrYRkQECRffwB9YB5a4gBHWmAxGQA7fmJaqH8UYAhO2wBRVvakLp7LPwCrCPS4LHVBODh%2BYZ4YKFp%2BGGg1uAB2uB5oKFX10%2FOoS6Obi%2FUgA}{verify part (a) of this exercise with Lean}.
\end{xca}

\subsection*{Hints for Lean}

\begin{leancode}
  import Mathlib.Algebra.Group.Subgroup.Basic

  variable {G : Type} [Group G] (H : Subgroup G)

  def Rel (H : Subgroup G) (g g' : G) : Prop := g⁻¹ * g' ∈ H

  lemma Rel_is_equivalence_relation : Equivalence (Rel H) := by
    constructor
    · sorry
    · sorry
    · sorry
\end{leancode}

A relation $R$ on a set $M$ is usually defined as a subset of $M \times M$.
In Lean, however, a relation on $M$ is defined as a function $M \to M \to \mathtt{Prop}$.
The idea is that for each $x, y \in M$ we obtain a proposition $P(x,y)$ which is {\ttfamily true} precisely when $(x,y) \in R$.
In our example, the function $\mathrm{Rel}$ sends $g \in G$ to the function that sends $g' \in G$ to the statement $g^{-1} \cdot g' \in H$.
You can read more about this on
\begin{center}
  \href{https://leanprover-community.github.io/logic_and_proof/relations_in_lean.html}{https://leanprover-community.github.io/logic\_and\_proof/relations\_in\_lean.html}.
\end{center}

The {\ttfamily constructor} tactic splits the current goal into three subgoals,
one for each defining property of an equivalence relation.
The bullet points are syntax to separate the three subgoals.

A \href{https://github.com/pedro-nunez/lean-algebra-2021/blob/master/equivalence-relation/solution.lean}{solution} is available in the repository
\begin{center}
  \href{https://github.com/pedro-nunez/lean-algebra-2021}{https://github.com/pedro-nunez/lean-algebra-2021}.
\end{center}

\subsection*{Lean's search tactics}

Many simple statements are already implemented in mathlib.
The {\ttfamily exact?} tactic searches for results already proven in mathlib
that can help solve the current goal.
For the first {\ttfamily sorry} above, try the following:

\begin{leancode}
  intro g
  unfold Rel
  simp only [inv_mul_cancel]
  exact?
\end{leancode}

Other search tactics you can try are {\ttfamily apply?}, {\ttfamily simp?} and {\ttfamily rw?}.

{\renewcommand{\thefootnote}{\fnsymbol{footnote}}\footnotetext[2]{Exercises designed for the lecture \href{https://home.mathematik.uni-freiburg.de/arithgeom/lehre/ws21/azt}{\emph{Algebra and Number Theory}}, for which I was a teaching assistant.
This was an introductory lecture on abstract algebra, polynomial rings and Galois theory, taught by Annette Huber-Klawitter at the University of Freiburg during the Winter Semester 2021/2022.}}

\end{document}
